% !TEX program = xelatex
% Documentation deck: every feature of the Sapienza Beamer theme, explained in the slides.
% Build with XeLaTeX:  make example   (light and dark variants)
\providecommand{\sapvariant}{light}   % light | dark
\documentclass[t,aspectratio=169]{beamer}
\usetheme[\sapvariant]{Sapienza}
\usepackage{booktabs}

\title[Sapienza Beamer Theme]{Sapienza Beamer theme}
\subtitle{Documentation and examples}
\author[Your Name]{Your Name - ID}
\institute{Sapienza University of Rome}
\date{}

% small helper to show source code in the slides (no verbatim, so frames stay non-fragile)
\newcommand{\code}[1]{{\ttfamily\small #1}}

\begin{document}

%% ------------------------------------------------------------------ title
\begin{frame}[plain]\titlepage\end{frame}

\begin{frame}{Credits}
  This theme is a LaTeX Beamer port of the \textbf{Sapienza PPT template} by Pietro Nardelli,
  inspired by the Sapienza NLP Group.
  \vspace{0.8em}
  \begin{itemize}
    \item \textbf{Original template:} \url{https://github.com/pietro-nardelli/sapienza-ppt-template}
    \item \textbf{Original author:} Pietro Nardelli
    \item \textbf{License:} {\itshape Attribution-NonCommercial-ShareAlike 4.0 International (CC BY-NC-SA 4.0)}
    \item \textbf{Note:} unofficial, not endorsed by Sapienza University of Rome
  \end{itemize}
\end{frame}

\begin{frame}[plain]
  \sapToc{Contents}{Getting started,Slide types,Content blocks,Custom components,Positioning,Reference}
\end{frame}

%% ------------------------------------------------------------------ 1. getting started
\begin{frame}{1. Getting started}
  \begin{itemize}
    \item Compile with \textbf{XeLaTeX} or LuaLaTeX: the theme uses \code{fontspec}
    \item Copy \code{beamerthemeSapienza.sty}, \code{fonts/} and \code{img/} next to your \code{.tex}
    \item Use \code{make} for both variants, or run \code{latexmk -xelatex}
    \item Run twice: the overlays use \code{remember picture} (latexmk does it for you)
  \end{itemize}
\end{frame}

\begin{frame}{Minimal document}
  \small
  \begin{itemize}
    \item \code{\textbackslash documentclass[t,aspectratio=169]\{beamer\}}
    \item \code{\textbackslash usetheme[light]\{Sapienza\}} \hfill (or \code{[dark]})
    \item \code{\textbackslash title[Short]\{Long title\}}, \code{\textbackslash author[Short]\{Full name\}}
    \item \code{\textbackslash subtitle}, \code{\textbackslash institute}
  \end{itemize}
  \vspace{0.6em}
  Use the class option \code{t}: it aligns frame content to the top, as in the original template.
  The short title and short author are printed in the footer.
\end{frame}

\begin{frame}{Light and dark variants}
  The variant is a theme option: \code{\textbackslash usetheme[dark]\{Sapienza\}}.
  Colours are defined once and switch with the option.
  \vspace{0.8em}
  \begin{itemize}
    \item \textbf{Accents} (same in both): teal \code{\#006778}, burgundy \code{\#6F0A19}
    \item \textbf{Light}: white background, grey text \code{\#595959}, banner \code{\#E9EDEE}
    \item \textbf{Dark}: background \code{\#353535}, text \code{\#E9EDEE}, banner \code{\#595959}
  \end{itemize}
  \vspace{0.6em}
  To pick the variant from the command line, define \code{\textbackslash sapvariant} before the
  document class, as this deck does (see the Makefile).
\end{frame}

%% ------------------------------------------------------------------ 2. slide types
\begin{frame}[plain]
  \sapSection{2. Slide types}{Five layouts, as in the PPT template}{%
    \textbf{Title slide}: \code{\textbackslash titlepage} in a \code{[plain]} frame.\par\medskip
    \textbf{Table of contents}: \code{\textbackslash sapToc}.\par\medskip
    \textbf{Content slide}: a normal \code{frame} with a title.\par\medskip
    \textbf{Section slide}: \code{\textbackslash sapSection}, the one you are looking at.\par\medskip
    \textbf{Closing slide}: \code{\textbackslash sapClosing}.}
\end{frame}

\begin{frame}{Title slide}
  \small
  \code{\textbackslash begin\{frame\}[plain]\textbackslash titlepage\textbackslash end\{frame\}}
  \vspace{0.6em}
  \begin{itemize}
    \item Title comes from \code{\textbackslash title}
    \item The three lines below are \code{\textbackslash author}, \code{\textbackslash subtitle} and \code{\textbackslash institute}
    \item Use \code{\textbackslash\textbackslash} inside \code{\textbackslash author} to add more lines
    \item The frame must be \code{[plain]}: no title bar, no footer
  \end{itemize}
\end{frame}

\begin{frame}{Table of contents}
  \small
  \code{\textbackslash sapToc\{Contents\}\{Section 1,Section 2,...\}}
  \vspace{0.6em}
  \begin{itemize}
    \item First argument: the slide title. Second: a comma-separated list of entries
    \item Entries are numbered automatically, two columns, row by row
    \item Up to six entries fit the original grid; use it inside a \code{[plain]} frame
    \item Separators alternate burgundy (left column) and teal (right column)
  \end{itemize}
\end{frame}

\begin{frame}{Content slide}
  A normal frame: \code{\textbackslash begin\{frame\}\{Title\}}. The title is 26\,pt bold, with the two
  accent bars above it. The body starts below the title and stops above the footer, which shows
  the short title, the short author and the slide number.
  \vspace{0.8em}
  \begin{itemize}
    \item The grey-blue banner on top and the burgundy corner with the logo are drawn by the theme
    \item Body text is 13\,pt, lists use grey dashes
    \item Body width matches the original placeholder (8.4 in on a 10 in slide)
  \end{itemize}
\end{frame}

%% ------------------------------------------------------------------ 3. content blocks
\begin{frame}{3. Lists}
  \begin{columns}[T]
    \begin{column}{0.48\textwidth}
      \textbf{Itemize}
      \begin{itemize}
        \item First level
          \begin{itemize}
            \item Second level
              \begin{itemize}\item Third level\end{itemize}
          \end{itemize}
        \item Another item
      \end{itemize}
    \end{column}
    \begin{column}{0.48\textwidth}
      \textbf{Enumerate}
      \begin{enumerate}
        \item First step
        \item Second step
        \item Third step
      \end{enumerate}
    \end{column}
  \end{columns}
\end{frame}

\begin{frame}{Text, links and math}
  Regular text, \textbf{bold}, \emph{emphasis} and a link: \url{https://github.com/pietro-nardelli/sapienza-ppt-template}.
  \vspace{0.8em}
  Inline math $e^{i\pi}+1=0$ and a displayed equation:
  \[
    \mathcal{L}(\theta)=-\frac{1}{N}\sum_{i=1}^{N}\log p_\theta(y_i\mid x_i)
  \]
  Links are blue (\code{\#1C3678}) in the light variant and cyan in the dark one.
\end{frame}

\begin{frame}{Tables}
  \centering
  \begin{tabular}{lcc}
    \toprule
    \textbf{Model} & \textbf{Accuracy} & \textbf{F1} \\
    \midrule
    Baseline & 71.2 & 69.8 \\
    Ours     & 78.5 & 77.1 \\
    \bottomrule
  \end{tabular}
  \vspace{1em}

  Use \code{booktabs} (load it in your preamble) for clean rules; colours follow the theme.
\end{frame}

\begin{frame}{Incremental reveal}
  Standard Beamer overlays work as usual:
  \begin{itemize}
    \item<1-> \code{\textbackslash item<1->} appears on the first click
    \item<2-> \code{\textbackslash item<2->} appears on the second
    \item<3-> and \code{\textbackslash pause}, \code{\textbackslash only}, \code{\textbackslash uncover} are available too
  \end{itemize}
  The slide number keeps counting frames, not overlays.
\end{frame}

%% ------------------------------------------------------------------ 4. custom components
\begin{frame}{4. Bad / Good comparison}
  \small
  \code{\textbackslash sapCompare\{\}\{y\}\{bad text\}\{good text\}}: \code{y} is the distance of the top of the boxes from the top
  of the slide, in inches of the original 10\,x\,5.625\,in slide.
  \sapCompare{}{3.4}%
    {A bad example: dense text, no structure, hard to read from the back of the room.}%
    {A good example: short sentences, one idea per slide, large and readable text.}
\end{frame}

\begin{frame}{Timeline}
  \sapAt[text width=8.4in]{0.798}{1.55}{\footnotesize\code{\textbackslash sapTimeline\{T1\}\{desc 1\}\{T2\}\{desc 2\}\{T3\}\{desc 3\}}: three milestones, the middle one hangs below the line.}
  \sapTimeline{Time 1}{First milestone of the project}%
              {Time 2}{Second milestone, described below the line in a wider box}%
              {Time 3}{Final milestone}
\end{frame}

\begin{frame}[plain]
  \sapSection{Section slide}{Subtitle goes here}{%
    \code{\textbackslash sapSection\{Title\}\{Subtitle\}\{Body\}} builds a split slide: title and
    subtitle on the left half, free body text on the right.\par\medskip
    It must be used inside a \code{[plain]} frame. The body is 11\,pt and may contain
    any text: lists, links or math.\par\medskip
    In the dark variant the right half is lighter grey.}
\end{frame}

%% ------------------------------------------------------------------ 5. positioning
\begin{frame}{5. Absolute positioning}
  \small
  \code{\textbackslash sapAt[options]\{x\}\{y\}\{text\}} puts a TikZ node with its top-left corner at $(x,y)$.
  Units are inches of the original slide (10\,x\,5.625), origin top-left, $y$ pointing down.
  \sapAt[text width=3cm,draw=sapTeal,inner sep=4pt]{1.0}{2.7}{Box at (1.0, 2.7) with a teal border.}
  \sapAt[text width=3.2cm,fill=sapBurg,text=white,inner sep=4pt]{4.5}{2.7}{Box at (4.5, 2.7) with a burgundy fill.}
  \sapAt[text width=3cm,text=sapGray]{7.5}{3.2}{Plain text at (7.5, 3.2).}
\end{frame}

\begin{frame}{Coordinates scale with the page}
  \begin{itemize}
    \item All positions are relative to the page width: \code{\textbackslash sapU} = 1 in of the original slide
    \item Font sizes use \code{\textbackslash sapfs\{size\}\{leading\}}, e.g.\ \code{\textbackslash sapfs\{26\}\{30\}} for a title
    \item So the theme works at any \code{aspectratio=169} paper size, not only 10 in wide
    \item For your own drawings use \code{sapcanvas}: a TikZ scope in the same units
  \end{itemize}
\end{frame}

%% ------------------------------------------------------------------ 6. reference
\begin{frame}{6. Colours and fonts}
  \begin{columns}[T]
    \begin{column}{0.5\textwidth}
      \textbf{Colours} (usable in any TikZ/xcolor command)
      \begin{itemize}
        \item \textcolor{sapTeal}{\code{sapTeal}}, \textcolor{sapBurg}{\code{sapBurg}}
        \item \code{sapGray}, \code{sapSilver}
        \item \code{sapBg}, \code{sapText}, \code{sapTitle}: they change with the variant
      \end{itemize}
    \end{column}
    \begin{column}{0.48\textwidth}
      \textbf{Fonts}
      \begin{itemize}
        \item Catamaran: titles and body
        \item {\sapRaleway Raleway}: table of contents (\code{\textbackslash sapRaleway})
        \item Files are loaded from \code{fonts/}
      \end{itemize}
    \end{column}
  \end{columns}
\end{frame}

\begin{frame}{Command summary}
  \small
  \begin{tabular}{@{}ll@{}}
    \code{\textbackslash titlepage} & title slide\\
    \code{\textbackslash sapToc\{title\}\{a,b,c\}} & table of contents\\
    \code{\textbackslash sapSection\{t\}\{sub\}\{body\}} & split section slide\\
    \code{\textbackslash sapCompare\{\}\{y\}\{bad\}\{good\}} & Bad / Good boxes\\
    \code{\textbackslash sapTimeline} & three-point timeline\\
    \code{\textbackslash sapClosing\{text\}} & closing slide\\
    \code{\textbackslash sapAt[opts]\{x\}\{y\}\{text\}} & absolute text box\\
    \code{sapcanvas} & TikZ scope in slide inches\\
  \end{tabular}
\end{frame}

%% ------------------------------------------------------------------ technical notes
\begin{frame}{Technical notes: files and units}
  \small
  \begin{itemize}
    \item \code{beamerthemeSapienza.sty}: the whole theme (colours, fonts, templates, components)
    \item \code{fonts/}: Catamaran (\code{.ttf}) and Raleway (\code{.otf}), loaded by \code{fontspec} with \code{Path=fonts/}
    \item \code{img/}: logos, four PNGs (\code{logo\_small|big\_light|dark})
    \item \code{\textbackslash sapU} = page width / 10: one inch of the original 10\,x\,5.625\,in slide
    \item \code{\textbackslash sapPt} = page width / 722: one point of the original slide
    \item \code{\textbackslash sapfs\{s\}\{l\}} sets font size and leading in \code{\textbackslash sapPt}
  \end{itemize}
\end{frame}

\begin{frame}{Technical notes: how a slide is drawn}
  \small
  \begin{itemize}
    \item \code{\textbackslash sapmode} (\code{content}, \code{title}, \code{toc}, \code{section}) is reset to \code{content} at each frame start
    \item \code{background canvas} paints the fills: body, grey banner, split halves of section slides
    \item \code{background} paints bars, corner block with logo, footer text and slide number
    \item Both use \code{sapcanvas}: a TikZ overlay with \code{remember picture}, origin at the page top-left
    \item This is why two compilation runs are needed: the first one stores the node positions
    \item The frame title is a fixed-height box, so the body always starts at the same height
  \end{itemize}
\end{frame}

\begin{frame}{Technical notes: customisation}
  \small
  \begin{itemize}
    \item Colours: change the \code{\textbackslash definecolor} lines of the \code{.sty}, one block per variant
    \item List markers: \code{\textbackslash setbeamertemplate\{itemize item\}} (grey dash, same for all levels)
    \item Margins: \code{text margin left/right} = 0.798\,in of the original slide
    \item Font sizes: \code{\textbackslash normalsize} 13\,pt, \code{\textbackslash small} 11\,pt, \code{\textbackslash footnotesize} 9\,pt
    \item Requirements: XeLaTeX or LuaLaTeX, \code{tikz}, \code{fontspec}, \code{xcolor}, \code{etoolbox}
    \item Aspect ratio: designed for \code{aspectratio=169}; other ratios are not supported
  \end{itemize}
\end{frame}

\begin{frame}{References}
  \begin{itemize}
    \item \textbf{Template:}
      \begin{itemize}\item \url{https://github.com/pietro-nardelli/sapienza-ppt-template}\end{itemize}
    \item \textbf{License:}
      \begin{itemize}\item {\itshape Attribution-NonCommercial-ShareAlike 4.0 International (CC BY-NC-SA 4.0)}\end{itemize}
    \item \textbf{Fonts:}
      \begin{itemize}\item Catamaran and Raleway, SIL Open Font License\end{itemize}
  \end{itemize}
\end{frame}

\begin{frame}[plain]
  \sapClosing{Thank you for the attention!}
\end{frame}

\end{document}
